\subsection{One chain interface}

\begin{definition}[Certified transfer chain]
\label{def:body-chain-notation}
\label{def:body-six-slot-chain}
A \emph{certified transfer chain} consists of consecutive V triangles
$T_{i_0},\ldots,T_{i_m}$, lower bounds
$x_0,\ldots,x_m$ for their incoming boundary reaches, and one justified
handoff at each intermediate edge.  Each handoff is one of the following.

\begin{enumerate}
\item An exact center-free or center-assisted transfer supplied by
Proposition~\ref{prop:body-transfer-interface}.
\item The identity relaxation $x_{j+1}\ge x_j$.  This relaxation is permitted
only when $T_{i_j}$ is nonsupercritical, the handoff edge is center-free, and
every nonincident role has zero-length trace on that edge.
\item A branch-specific affine or threshold map that is nondecreasing and
proved pointwise below the exact transfer on its stated input interval, or a
quarter-radial or Vd-corner lower bound.  The domains and pointwise statements
are recorded in Appendix~\ref{sec:strategy2-optimization-problems}.
\end{enumerate}

The order of the chain is geometric: the first recorded handoff acts first.
The only functional relaxations used are the identity, one affine chord on a
selected $T_+$ branch, and one low-root threshold; their formulas and validity
domains are Definitions~\ref{def:s2-pure-relaxations}
and~\ref{def:s2-pure-strict-certificate}.  Lemma~\ref{lem:reader-relaxed-composition}
shows that replacing exact nondecreasing transfer maps by these certified
pointwise lower bounds preserves the returned lower bound.
Writing an identity slot asserts only the individual inequality
$x_{j+1}\ge x_j$; it does not assert a separate pointwise comparison between
two formal compositions.
\end{definition}

\begin{definition}[Center-free ordinary boundary path]
\label{def:body-center-free-path}
A consecutive boundary path is \emph{center-free ordinary} when, on every
internal edge, the center role has empty trace and every role other than the
two incident path roles has zero-length trace.  Endpoint contributions from
roles outside the path must be incorporated into the two endpoint residuals
before the path is used.
\end{definition}

\begin{proposition}[Four possible endings of a transfer chain]
\label{prop:body-strategy2-chain-principle}
Every Strategy~2 certificate has one of the following endings.

\begin{enumerate}
\item \emph{Strict cycle:} six center-free nonsupercritical handoffs force
$A_0<A_1<\cdots<A_5<A_0$.
\item \emph{Endpoint deficit:} after all external endpoint contributions are
absorbed, the two endpoint capacities have sum strictly below one; the
center-free path budget is then impossible.
\item \emph{Returned-demand excess:} a five-V-triangle chain returns a lower
bound $Z$ while the remaining role has the strict upper bound
$A_0<1-H$, and the analytic certificate gives $Z>1-H$.
\item \emph{Radial separation:} the analytic endpoint bounds leave a
nonempty interval between the vertex-side and center-side radial traces.
\end{enumerate}
\end{proposition}

\begin{proof}
Proposition~\ref{prop:body-transfer-interface} gives the handoff induction,
and Lemma~\ref{lem:reader-boundary-path} gives the endpoint conclusion.  The
exact one-gap interface gives the returned-demand upper bound.  The
branch-specific statements in
Proposition~\ref{prop:body-special-terminal-certificates} handle the radial
cases.
\end{proof}
